% !TEX program = xelatex
\documentclass[11pt]{article}
\usepackage[margin=1.1in]{geometry}
\usepackage{fontspec}
\setmainfont{DejaVu Serif}[Scale=0.9]
\usepackage{polyglossia}
\setmainlanguage{english}
\usepackage{amsmath,amssymb,amsthm,mathtools}
\usepackage{booktabs,array,enumitem}
\usepackage{xcolor}
\usepackage[colorlinks=true,linkcolor=blue!50!black,citecolor=blue!50!black,urlcolor=blue!50!black]{hyperref}
\usepackage{attachfile2}

\newtheorem{theorem}{Theorem}[section]
\newtheorem{proposition}[theorem]{Proposition}
\newtheorem{lemma}[theorem]{Lemma}
\newtheorem{corollary}[theorem]{Corollary}
\theoremstyle{definition}
\newtheorem{conjecture}[theorem]{Conjecture}
\newtheorem{question}{Question}
\theoremstyle{remark}
\newtheorem{remark}[theorem]{Remark}

\newcommand{\tagP}{\textsf{\small[proved here]}}
\newcommand{\tagC}{\textsf{\small[cited]}}
\newcommand{\tagN}{\textsf{\small[numerically checked]}}
\newcommand{\tagK}{\textsf{\small[kernel-checked arithmetic]}}
\newcommand{\tagO}{\textsf{\small\textcolor{red!70!black}{[open]}}}

\newcommand{\C}{\mathbb{C}}\newcommand{\R}{\mathbb{R}}\newcommand{\Z}{\mathbb{Z}}
\newcommand{\M}{\mathfrak{M}}\newcommand{\Rep}{\mathrm{Rep}}\newcommand{\SRep}{\mathrm{SRep}}
\newcommand{\BD}{\mathrm{BD}}
\newcommand{\bundleurl}{https://totogt.github.io/geometry/collab/pinot-orthogonal-quivers/}

\title{Affine orthogonal quiver varieties at $\delta$ and $2\delta$,\\ and their contact links}
\author{Pablo Nogueira Grossi\\[2pt]
{\small G6 LLC (d/b/a PrincipiaOrthogona), Newark, NJ, USA \quad \texttt{g6llc@proton.me}\quad ORCID 0009-0000-6496-2186}}
\date{Version 0.2 --- 5 October 2026 --- prepared for evaluation by V.~Pinot}

\begin{document}
\maketitle

\begin{abstract}
Pinot recently computed the quiver varieties $\M_\delta$ and $\M_{2\delta}$ of the orthogonal quivers whose underlying graph is an affine Dynkin diagram, and identified Kleinian surface singularities among them \cite{Pinot}. We make three observations and one construction. (i) Using only the closed immersion of Li \cite[Prop.~9.2.1]{Li19}, whenever $\M_\delta$ is a surface the involution acts trivially and $\M_\delta$ is the classical Kleinian surface of the underlying affine diagram. (ii) Consequently the two $\widetilde D$ rows of \cite[Thm.~1.2]{Pinot} are of type $D_{2n-2}$ and $D_{2n-1}$, not $D_{n+1}$; for $(\widetilde D_{2n-1},a)$ the defining relation is $z^2=z(-x)^{n-1}-xy^2$. We confirm both by an independent numerical computation of the invariant ring. (iii) In every row at $2\delta$ the group $\Gamma$ contains the classical group of the underlying graph with index $2$ (or $1$); we conjecture that the $2\delta$ surfaces are quotients of the classical surface by an involution, and prove the corresponding statement for the coordinate rings. Finally, each surface is a cone whose link is a contact $3$-manifold $S^3/\Gamma$; we show that the symmetric representation space is quaternionic, so the link is the $3$-Sasakian (Kempf--Ness) quotient of the unit sphere, and we tabulate the Reeb (Hopf) Seifert fibrations.
\end{abstract}

\begin{center}\fbox{\parbox{0.93\textwidth}{\small
\textbf{How to read this version.} Every statement carries one tag: \tagP, \tagC, \tagN{} (a script in the bundle; evidence, not proof), \tagK{} (Lean~4, arithmetic only), or \tagO. \textbf{Verification bundle} (Python and Lean, with a README): \textattachfile[color=0 0 0.6]{DM3QUIVERS-verification-bundle.zip}{\textcolor{blue!50!black}{\underline{click here to open the attached zip}}} (attached to this PDF), or download it from \url{\bundleurl}. This text was prepared with the assistance of Claude (Anthropic); the author takes responsibility for it, and it has not yet been checked by V.~Pinot.}}\end{center}

\section{Introduction}

Nakajima quiver varieties of affine Dynkin type recover Kronheimer's ALE spaces and the Kleinian singularities $\C^2/\Gamma$ \cite{Kron,Nak94}; the ADE labels go back to Gabriel's classification of quivers of finite type \cite{Gab72}. Pinot studies the analogous quotients for \emph{orthogonal quivers} \cite{Pinot}: quivers $Q$ with a contravariant involution $\tau$ and a sign $s$, whose symmetric representations $\SRep(\overline Q,\mathrm d)$ carry an action of $G^b_{\mathrm d}=\prod O(\mathrm d_i)\times\prod GL(\mathrm d_j)$, and
\[\M_{\mathrm d}=\mu^{-1}(0)/\!/G^b_{\mathrm d}.\]
His Theorem~1.2 lists $\M_\delta$ and $\M_{2\delta}$ for ten infinite families. This note re-reads that table from two angles: the comparison with the classical quiver variety of the same graph (\S\S3--4), and the contact geometry of the link of each surface (\S5).

\paragraph{Results.}
\begin{itemize}[leftmargin=*]
\item \textbf{Theorem A} (\S3). If $\M_\delta$ is a surface, then $\M_\delta\cong\M_\delta^{\rm cl}$, the classical Kleinian surface of the underlying affine graph, and the involution acts trivially on it. Conditional only on the applicability of Li's closed immersion to Pinot's unframed setting.
\item \textbf{Corollary B} (\S3). The $\widetilde D$ rows read $D_{2n-2}$ and $D_{2n-1}$. The $(\widetilde D_{2n-2},v)$ equation of \cite{Pinot} is correct; the $(\widetilde D_{2n-1},a)$ relation is $z^2=z(-x)^{n-1}-xy^2$. Confirmed numerically for $n=3,4,5$ by computing the invariant ring directly.
\item \textbf{Conjecture C} (\S4). Each $2\delta$ surface is $\M^{\rm cl}_\delta/\iota$ for an involution $\iota$; equivalently $\Gamma_{2\delta}\supset\Gamma_{\rm cl}$ with index $\le2$, and the classical link double-covers the $2\delta$ link. The ring-level statement is proved.
\item \textbf{Theorem D} (\S5). $\SRep(\overline Q,\mathrm d)$ is a quaternionic subspace; the link is the Kempf--Ness / $3$-Sasakian quotient of the unit sphere; the Reeb flow is the Hopf flow, with the Seifert data of Table~\ref{tab:seifert}.
\end{itemize}

\section{Setting}\label{sec:setup}
We use the notation of \cite{Pinot} throughout, with the standard signed space ($J_i=\mathrm{id}$ in suitable bases), so that the involution on representations is $\sigma(v)_a=s(a)\,v_{\tau a}^{T}$ and $\SRep=\Rep^\sigma$ \cite[Def.~2.3, 2.5]{Pinot}. The moment map is the restriction of the quiver moment map, $\mu(\SRep)\subset\mathfrak g^b$ \cite[Cor.~2.13]{Pinot}. The invariant ring is generated by traces of cycles \cite[Thm.~3.9]{Pinot}, \cite{LP90}. The scaling $t\cdot v$ gives $\operatorname{tr}(\Lambda^kC)$ the weight $k\,\ell(C)$, so every $\M_{\mathrm d}$ is a cone \cite[\S3.5.1]{Pinot}.

\paragraph{The comparison map.} Let $\M^{\rm cl}_{\mathrm d}=\mu^{-1}(0)/\!/G_{\mathrm d}$ be the classical quiver variety of the same quiver. By universality of categorical quotients there is a morphism $\iota:\M_{\mathrm d}\to(\M^{\rm cl}_{\mathrm d})^{[\sigma]}$. Remark~2.15 of \cite{Pinot} states that it is an isomorphism, citing \cite{Li19,Nak25}. The sources give less: Li proves that $\iota$ is a \emph{closed immersion} for an arbitrary graph \cite[Prop.~9.2.1]{Li19}, calls the isomorphism conjectural, and proves it in type $A_n$ with alternating signs \cite[Prop.~9.2.2]{Li19}; Nakajima treats only smooth $\sigma$-quiver varieties (generic $\zeta$) \cite{Nak25}. Li works with framed varieties $(v,w)$ and prescribed forms; we \emph{assume} throughout that his Proposition~9.2.1 applies to Pinot's unframed signed setting (Question~\ref{q:li}). We use only the closed immersion.

\section{The surfaces at $\delta$}\label{sec:delta}

\begin{theorem}\label{thm:A}
Assume Li's closed immersion $\iota:\M_\delta\hookrightarrow\M^{\rm cl}_\delta$ holds in Pinot's setting, $\delta$ being the minimal imaginary root of the underlying affine graph. If $\M_\delta$ is a reduced irreducible surface (as in all non-trivial rows of \cite[Thm.~1.2]{Pinot}), then $\iota$ is an isomorphism and $[\sigma]$ acts trivially on $\M^{\rm cl}_\delta$. In particular $\M_\delta\cong\C^2/\Gamma$ with $\Gamma$ the McKay group of the underlying affine graph.
\end{theorem}
\begin{proof}
$\M^{\rm cl}_\delta\cong\C^2/\Gamma$ is an irreducible surface \cite{Kron,Nak94}. A closed immersion of a reduced irreducible surface into an irreducible surface is surjective, hence an isomorphism onto it. Since its image lies in the fixed locus of $[\sigma]$, that fixed locus is everything. \tagP{} (conditional on the hypothesis).
\end{proof}

\begin{corollary}\label{cor:B}
The non-trivial $\delta$ rows are of type $A_{2n-2}$ for $\widetilde A_{2n-2}$, $A_{2n-1}$ for $\widetilde A_{2n-1}$, $D_{2n-2}$ for $(\widetilde D_{2n-2},v)$ and $D_{2n-1}$ for $(\widetilde D_{2n-1},a)$. The last two replace the label $D_{n+1}$ printed in \cite[Thm.~1.2]{Pinot}; they agree with it only for $n=3$ in the first case, and never in the second.
\end{corollary}

Corollary~\ref{cor:B} depends on the hypothesis of Theorem~\ref{thm:A}. The next two results check it independently.

\begin{proposition}[weighted monomials]\label{prop:enum}
In the $\widetilde D$ rows let $x=\det C_3$, $y=\operatorname{tr}(C_1C_2)$, $z=\operatorname{tr}(C_1C_2C_3)$, with $L=2n-5$ for $(\widetilde D_{2n-2},v)$ and $L=2n-4$ for $(\widetilde D_{2n-1},a)$ \cite[p.~18]{Pinot}. Their weights are $(4,2L+2,2L+4)$, and a relation homogeneous with $z^2$ has degree $4L+8$. The only monomials of that degree are:
\begin{itemize}[leftmargin=*]
\item $(\widetilde D_{2n-2},v)$, $n\ge4$: $z^2,\ xy^2,\ x^{n-1}y,\ x^{2n-3}$. The relation $z^2+\beta xy^2+\alpha x^{n-1}y+\gamma x^{2n-3}$ is of type $D_{2n-2}$ if and only if $\beta\ne0$ and $\gamma\neq\alpha^2/4\beta$. For $n=3$ the weights are $(4,4,6)$, so $y^3$ also has degree $12$: the relation is $z^2+c(x,y)$ with $c$ a binary cubic, of type $D_4$ if and only if $c$ has three distinct linear factors.
\item $(\widetilde D_{2n-1},a)$: $z^2,\ x^{n-1}z,\ xy^2,\ x^{2n-2}$. The relation $z^2+\alpha x^{n-1}z+\beta xy^2+\gamma x^{2n-2}$ is of type $D_{2n-1}$ if and only if $\beta\ne0$ and $\gamma\ne\alpha^2/4$. The term $x^{n-1}y$ printed in \cite[Lemma~3.27]{Pinot} has degree $4L+6$ and cannot occur.
\end{itemize}
\end{proposition}
\begin{proof}
Solve $4a+(2L+2)b+(2L+4)c=4L+8$ in non-negative integers. Then complete the square: in $y$ for the first row, in $z$ for the second. The result is the normal form $z^2+xy^2+c\,x^{m-1}$ of $D_m$ with $m=2n-2$, resp.\ $m=2n-1$, and $c\neq0$ exactly under the stated conditions. Equivalently, by the Milnor--Orlik formula $\mu=\prod(d/w_i-1)=L+3$ \cite{MO}. \tagP{} \tagK
\end{proof}

\begin{proposition}[direct computation]\label{prop:numeric}
For $n=3,4,5$ we sampled $80$ random points of $\mu^{-1}(0)\cap\SRep(\overline Q,\delta)$ by Gauss--Newton and found all linear relations among the monomials of weight $\le 4L+8$ in $x,y,z$. In each case there is exactly one relation (relative singular value $\le2\cdot10^{-8}$, against $\ge0.1$ for the next), with integer coefficients:
\[
(\widetilde D_{2n-2},v):\ z^2=-y(-x)^{n-1}-xy^2,\qquad (\widetilde D_{2n-1},a):\ z^2=z(-x)^{n-1}-xy^2 .
\]
The first is the equation of \cite[Lemma~3.28]{Pinot}. The second is \cite[Lemma~3.27]{Pinot} with $y$ replaced by $z$ in the first term: the step $\operatorname{tr}(C_1C_2^2C_3^3)=\operatorname{tr}(C_1C_2C_3)(-\det C_3)^{n-1}$ appears to have been written with $\operatorname{tr}(C_1C_2)$. In both rows the non-degeneracy conditions of Proposition~\ref{prop:enum} hold ($\beta=1$, discriminant $-1/4$; for $n=3$ in the first row the cubic is $xy(x+y)$, with three distinct factors). For sign $-1$ in $(\widetilde D_{2n-1},a)$, all invariants vanish, in agreement with \cite{Pinot}. \tagN
\end{proposition}

\begin{table}[h]\centering\small
\begin{tabular}{@{}lllllll@{}}
\toprule
quiver & sign & d & equation & type & $\Gamma$ & $|\Gamma|$\\\midrule
$(\widetilde A_{2n-2},v\text{-}a)$ & $+1$ & $\delta$ & $xy=z^{2n-1}$ & $A_{2n-2}$ & $\Z_{2n-1}$ & $2n-1$\\
 & $-1$ & $2\delta$ & $xy=z^{4n-2}$ & $A_{4n-3}$ & $\Z_{4n-2}$ & $4n-2$\\
$(\widetilde A_{2n-1},a\text{-}a)$ & $(+,+)$ & $\delta$ & $xy=z^{2n}$ & $A_{2n-1}$ & $\Z_{2n}$ & $2n$\\
 & $(+,-)$ & $2\delta$ & $xy=z^{4n}$ & $A_{4n-1}$ & $\Z_{4n}$ & $4n$\\
 & $(-,-)$ & $2\delta$ & $xy=z^{2n}$ & $A_{2n-1}$ & $\Z_{2n}$ & $2n$\\
$(\widetilde A_{2n-1},v\text{-}v)$ & & $\delta$ & $xy=z^{2n}$ & $A_{2n-1}$ & $\Z_{2n}$ & $2n$\\
$(\widetilde A_{2n-1},c)$ & & $2\delta$ & $x^{n+1}-y^2x-z^2=0$ & $D_{n+2}$ & $\BD_{4n}$ & $4n$\\
$(\widetilde D_{2n-2},v)$ & & $\delta$ & $z^2=-y(-x)^{n-1}-xy^2$ & $\mathbf{D_{2n-2}}$ & $\BD_{8(n-2)}$ & $8n-16$\\
$(\widetilde D_{2n-1},a)$ & $+1$ & $\delta$ & $\mathbf{z^2=z(-x)^{n-1}-xy^2}$ & $\mathbf{D_{2n-1}}$ & $\BD_{4(2n-3)}$ & $8n-12$\\\bottomrule
\end{tabular}
\caption{The non-trivial rows of \cite[Thm.~1.2]{Pinot}, with the corrections in bold. $\BD_{4k}$ is the binary dihedral group of order $4k$ (for $n=1$ in the $(\widetilde A_{2n-1},c)$ row, $D_3=A_3$). All other rows are reduced points.}\label{tab:main}
\end{table}

\section{The surfaces at $2\delta$}\label{sec:2delta}
At $2\delta$ the classical variety is $4$-dimensional, and the orthogonal surfaces are new. Comparing each $\Gamma_{2\delta}$ with the classical group $\Gamma_{\rm cl}$ of the same graph:
\begin{center}\small
\begin{tabular}{@{}llll@{}}\toprule
$2\delta$ row & $\Gamma_{2\delta}$ & $\Gamma_{\rm cl}$ & index\\\midrule
$(\widetilde A_{2n-2},v\text{-}a)$, $-1$ & $\Z_{4n-2}$ & $\Z_{2n-1}$ & 2\\
$(\widetilde A_{2n-1},a\text{-}a)$, $(+,-)$ & $\Z_{4n}$ & $\Z_{2n}$ & 2\\
$(\widetilde A_{2n-1},c)$ & $\BD_{4n}$ & $\Z_{2n}$ & 2\\
$(\widetilde A_{2n-1},a\text{-}a)$, $(-,-)$ & $\Z_{2n}$ & $\Z_{2n}$ & 1\\\bottomrule
\end{tabular}
\end{center}
The two index-$2$ overgroups of $\Z_{2n}$ in $SU(2)$ that occur here, $\Z_{4n}$ and $\BD_{4n}$, both appear.

\begin{conjecture}\label{conj:C}
Each $2\delta$ surface is isomorphic to $\M^{\rm cl}_\delta/\iota$ for an involution $\iota$ of the classical surface, trivial exactly in the $(-,-)$ row.
\end{conjecture}

\begin{proposition}\label{prop:rings}
Write $\C^2/\Z_m=\{XY=Z^m\}$ with $X=u^m$, $Y=v^m$, $Z=uv$.
\begin{enumerate}[label=(\alph*)]
\item The generator of $\Z_{2m}\supset\Z_m$ induces $\iota_1:(X,Y,Z)\mapsto(-X,-Y,Z)$. Its invariants $X^2,Y^2,Z$ satisfy $X^2Y^2=Z^{2m}$, the surface $A_{2m-1}$.
\item For $m=2n$, $j\in\BD_{4n}$ induces $\iota_2:(X,Y,Z)\mapsto(Y,X,-Z)$. Its invariants $s=X+Y$, $w=Z^2$, $t=Z(X-Y)$ satisfy $t^2=ws^2-4w^{n+1}$, the surface $D_{n+2}$.
\end{enumerate}
With $m=2n-1$, resp.\ $m=2n$, these reproduce the three index-$2$ rows above exactly.
\end{proposition}
\begin{proof}
Direct substitution; the invariant rings are the classical Klein invariants of $\Z_{2m}$ and $\BD_{4n}$. \tagP{} \tagN{} (symbolic, $m\le11$).
\end{proof}

Pinot's own generators fit this picture. In the $(+,-)$ row the $2\delta$ coordinates are $\det C_1,\det C_2$ and $\operatorname{tr}C_3$, where the classical $\delta$ coordinates are $\operatorname{tr}C_1,\operatorname{tr}C_2,\operatorname{tr}C_3$. The weights double for the first two, just as $X\mapsto X^2$, $Y\mapsto Y^2$ in~(a). The remaining step is to construct $\iota$ from the quiver data. \tagO

\section{The contact links}\label{sec:links}

\subsection{Cones and weights}
\begin{lemma}\label{lem:weights}
In every row of Table~\ref{tab:main}, Pinot's path-length weights of the three generators equal the degrees of the corresponding Klein invariants in $\C[u,v]^\Gamma$: $(m,m,2)$ for $\Z_m$ and $(4,2k,2k+2)$ for $\BD_{4k}$. Hence the identification $\M_{\mathrm d}\cong\C^2/\Gamma$ can be chosen graded, and the unit circle of the scaling $\C^*$ acts on $\C^2/\Gamma$ by scalars.
\end{lemma}
\begin{proof}
At $\delta$ for $\widetilde A_{m-1}$: $\operatorname{tr}C_1,\operatorname{tr}C_2$ have weight $m$, $\operatorname{tr}C_3$ weight $2$. At $2\delta$, determinants double the cycle length: $(4n-2,4n-2,2)$ and $(4n,4n,2)$. In the $(\widetilde A_{2n-1},c)$ row the weights are $(4,2n,2n+2)$, and in the $\widetilde D$ rows $(4,2L+2,2L+4)$ with $k=L+1$. The normal-form changes of variables used above preserve the grading. \tagP
\end{proof}

The link $L_Q=(\M_{\mathrm d}\setminus0)/\R_{>0}\cong S^3/\Gamma$ carries its canonical contact structure: the complex tangencies, unique up to contactomorphism \cite{CNPP}, which is the quotient of the standard structure of $S^3$ since $\Gamma\subset SU(2)$. It is universally tight and Milnor fillable. \tagC

\subsection{A real structure}
A real contact form needs a real exact symplectic form. Pinot's $\omega$ and $\mu$ are holomorphic, so we use the flat hyperk\"ahler structure of $T^*\Rep(Q,\mathrm d)$: $I=i$ and $J(x,y)=(-y^\dagger,x^\dagger)$ arrow by arrow.

\begin{lemma}\label{lem:quat}
In the standard signed space, $\sigma\circ J=J\circ\sigma$ and $\sigma$ is unitary. Hence $\SRep(\overline Q,\mathrm d)$ is a quaternionic subspace. The compact group $K^b=G^b\cap\prod U(\mathrm d_i)=\prod O(\mathrm d_i,\R)\times\prod U(\mathrm d_j)$ preserves $I,J,K$, and both moment maps send $\SRep$ into $\mathfrak k^b$.
\end{lemma}
\begin{proof}
$\sigma(Jv)_a=s(a)(-v_{(\tau a)^*}^\dagger)^T=-s(a)\overline{v_{(\tau a)^*}}=J(\sigma v)_a$, and similarly for $a^*$, using $s(a^*)=s(a)$ and $\tau(a^*)=\tau(a)^*$. $\sigma$ permutes matrix entries up to sign and transposition, so it preserves the Hermitian norm. For $\mu_\C$ the last claim is \cite[Prop.~2.11]{Pinot}; for $\mu_\R$ the same computation applies. \tagP{} \tagN
\end{proof}

\begin{theorem}\label{thm:D}
There is a $\C^*$-equivariant homeomorphism $\bigl(\mu_\R^{-1}(0)\cap\mu_\C^{-1}(0)\cap\SRep\bigr)/K^b\cong\M_{\mathrm d}$. Intersecting with the unit sphere gives
\[L_Q\cong\bigl(\mu_\R^{-1}(0)\cap\mu_\C^{-1}(0)\cap S(\SRep)\bigr)/K^b,\]
the $3$-Sasakian quotient of the unit sphere in the quaternionic space $\SRep$ \cite{BGM,BG98}. The smooth structure of $L_Q$ comes from Lemma~\ref{lem:weights}. The reduction need not be locally free on the zero set.
\end{theorem}
\begin{proof}
$\mu_\C^{-1}(0)\cap\SRep$ is a closed $G^b$-stable subvariety of a unitary representation of $G^b=(K^b)_\C$. Kempf--Ness \cite{KN} identifies its categorical quotient with the $K^b$-quotient of the zero set of the $K^b$-moment map, which is $\mu_\R$ by Lemma~\ref{lem:quat}. Both sides are cones, so the identification restricts to links. \tagP{} (standard)
\end{proof}

\subsection{Reeb flow}
By Lemma~\ref{lem:weights}, the Reeb flow of the contact form induced by the flat metric is the Hopf flow $x\mapsto e^{i\theta}x$ on $S^3/\Gamma$.

\begin{proposition}\label{prop:seifert}
The Hopf flow on $S^3/\Gamma$ is a Seifert fibration with the following data. The generic orbit has period $2\pi$ if $-1\notin\Gamma$ and $\pi$ otherwise. The exceptional fibre over a cone point of order $a$ has period (generic)$/a$.
\end{proposition}

\begin{table}[h]\centering\small
\begin{tabular}{@{}lllll@{}}\toprule
type & $\Gamma$ & $L_Q$ & base orbifold & generic period\\\midrule
$A_{m-1}$, $m$ odd & $\Z_m$ & $L(m,m-1)$ & $S^2(m,m)$ & $2\pi$\\
$A_{m-1}$, $m$ even & $\Z_m$ & $L(m,m-1)$ & $S^2(m/2,m/2)$ & $\pi$\\
$D_{k+2}$ & $\BD_{4k}$ & prism manifold & $S^2(2,2,k)$ & $\pi$\\\bottomrule
\end{tabular}
\caption{Seifert data of the Reeb flow: base orbifold and multiplicities. The Seifert invariants $\beta_i$ and the Euler number are not computed here.}\label{tab:seifert}
\end{table}
\begin{proof}
The stabiliser of $x$ is $\{\theta: e^{i\theta}x\in\Gamma x\}$. For generic $x$ only the scalars $\pm1\in\Gamma$ contribute. The base is $\C P^1/\bar\Gamma$, with $\bar\Gamma\subset SO(3)$ cyclic of order $m$ or $m/2$, or dihedral of order $2k$. Its cone points are the points of $\C P^1$ with non-trivial stabiliser in $\bar\Gamma$. The identity $\chi^{\rm orb}=2/|\bar\Gamma|$ checks the data. \tagP{} \tagN{} \tagK
\end{proof}

Under Conjecture~\ref{conj:C}, $S^3/\Gamma_{\rm cl}\to S^3/\Gamma_{2\delta}$ is a double covering of contact manifolds commuting with the Reeb flows: the link at $2\delta$ is a $\Z_2$-quotient of the classical link.

\subsection{Fillings}
The classical minimal resolution of $\C^2/\Gamma$ (Kronheimer's $\zeta\ne0$ quotient) is a strong symplectic (K\"ahler) filling of $L_Q$. It is not Stein, since it contains the exceptional $(-2)$-curves. For ADE singularities it is diffeomorphic to the Milnor fibre \cite{Brieskorn}, which is a Stein filling. For lens spaces with this contact structure all fillings are classified by Lisca \cite{Lis}. Whether Pinot's orthogonal quotient at $\zeta\ne0$ gives a resolution is open: only the $GL$ factors of $G^b$ admit characters, and the involution sends $\zeta$ to $-\zeta$ before a Weyl correction \cite{Nak25}. \tagC{} \tagO

\section{Questions for V.~Pinot}
\begin{question} Do you confirm the labels $D_{2n-2}$, $D_{2n-1}$ and the relation $z^2=z(-x)^{n-1}-xy^2$ of Corollary~\ref{cor:B} and Proposition~\ref{prop:numeric}?\end{question}
\begin{question}\label{q:li} Does Li's Proposition~9.2.1 apply to your unframed signed setting (for instance with $w=0$, or with an auxiliary framing vertex)? Theorem~\ref{thm:A} uses only this closed immersion.\end{question}
\begin{question} Do you have a source for the isomorphism of Remark~2.15 at $\zeta=0$ beyond type $A_n$ with alternating signs?\end{question}
\begin{question} Do the $2\delta$ surfaces arise as quotients of the classical surface by an involution (Conjecture~\ref{conj:C})? Is there a natural candidate for $\iota$ in terms of the quiver?\end{question}
\begin{question} Do you plan to treat the case of your Remark~1.3, $(\widetilde D_{2n-1},a)$ with sign $-1$ at $2\delta$? Under Conjecture~\ref{conj:C} one would expect an overgroup of $\BD_{4(2n-3)}$ of index at most~$2$.\end{question}

\appendix
\section{Verification bundle}
\textattachfile[color=0 0 0.6]{DM3QUIVERS-verification-bundle.zip}{\textcolor{blue!50!black}{\underline{Attached zip}}}; also at \url{\bundleurl}. Python~3 with \texttt{numpy}, \texttt{scipy}, \texttt{sympy}; Lean~4 core (no Mathlib).
\begin{itemize}[leftmargin=*]
\item \texttt{verify\_Dtilde\_relation\_numeric.py}: Proposition~\ref{prop:numeric}; saved output included.
\item \texttt{verify\_ade\_types.py}: Milnor numbers of the $\widetilde D$-row polynomials; homogeneity defect of the printed relation.
\item \texttt{verify\_2delta\_index2.py}: Proposition~\ref{prop:rings}.
\item \texttt{verify\_quaternionic\_SRep.py}: Lemma~\ref{lem:quat}.
\item \texttt{verify\_seifert.py}: Table~\ref{tab:seifert} and the group orders.
\item \texttt{DM3Quivers\_Check.lean}: weight degrees, the Milnor--Orlik numerator identity $(d-w_x)(d-w_y)(d-w_z)=(L+3)w_xw_yw_z$, group orders, and the orbifold identities. It does not check Remark~2.15 or any geometric statement.
\end{itemize}

\begin{thebibliography}{99}
\bibitem{BG98} C.~P.~Boyer, K.~Galicki, \emph{3-Sasakian manifolds}, Surveys in Differential Geometry~VI (1999); arXiv:hep-th/9810250.
\bibitem{BGM} C.~P.~Boyer, K.~Galicki, B.~M.~Mann, \emph{The geometry and topology of 3-Sasakian manifolds}, J.\ Reine Angew.\ Math.\ 455 (1994), 183--220.
\bibitem{Brieskorn} E.~Brieskorn, \emph{Singular elements of semi-simple algebraic groups}, Actes du Congr\`es International des Math\'ematiciens (Nice, 1970), Tome~2, 279--284.
\bibitem{CNPP} C.~Caubel, A.~N\'emethi, P.~Popescu-Pampu, \emph{Milnor open books and Milnor fillable contact 3-manifolds}, Topology 45 (2006), 673--689.
\bibitem{Gab72} P.~Gabriel, \emph{Unzerlegbare Darstellungen I}, Manuscripta Math.\ 6 (1972), 71--103.
\bibitem{KN} G.~Kempf, L.~Ness, \emph{The length of vectors in representation spaces}, Algebraic Geometry (Copenhagen 1978), Lecture Notes in Math.\ 732, Springer, 1979, 233--243.
\bibitem{Kron} P.~B.~Kronheimer, \emph{The construction of ALE spaces as hyper-K\"ahler quotients}, J.\ Differential Geom.\ 29 (1989), 665--683.
\bibitem{LP90} L.~Le Bruyn, C.~Procesi, \emph{Semisimple representations of quivers}, Trans.\ Amer.\ Math.\ Soc.\ 317 (1990), 585--598.
\bibitem{Li19} Y.~Li, \emph{Quiver varieties and symmetric pairs}, Represent.\ Theory 23 (2019), 1--56.
\bibitem{Lis} P.~Lisca, \emph{On symplectic fillings of lens spaces}, Trans.\ Amer.\ Math.\ Soc.\ 360 (2008), 765--799.
\bibitem{MO} J.~Milnor, P.~Orlik, \emph{Isolated singularities defined by weighted homogeneous polynomials}, Topology 9 (1970), 385--393.
\bibitem{Nak94} H.~Nakajima, \emph{Instantons on ALE spaces, quiver varieties, and Kac--Moody algebras}, Duke Math.\ J.\ 76 (1994), 365--416.
\bibitem{Nak25} H.~Nakajima, \emph{Instantons on ALE spaces for classical groups, involutions on quiver varieties, and quantum symmetric pairs}, arXiv:2510.13007 (2025).
\bibitem{Pinot} V.~Pinot, \emph{Quiver varieties for affine orthogonal quivers}, arXiv:2609.39434 (2026).
\end{thebibliography}
\end{document}
